\documentclass[11pt]{article}

\usepackage[margin=1in]{geometry}
\usepackage{amsmath,amssymb,amsthm}
\usepackage{graphicx}
\usepackage[hidelinks]{hyperref}

\newtheorem{theorem}{Theorem}
\newtheorem{lemma}{Lemma}
\newtheorem{corollary}{Corollary}
\newtheorem{remark}{Remark}

\newcommand{\CWOS}{\textit{CWO-S}}
\newcommand{\eps}{\varepsilon}
\newcommand{\one}{\mathbf 1}

\title{Strictly Convex Targets in the \(L_1(\mu,X)\) \(\CWOS\) Question}
\author{Run fa\_banach\_001, lane 7}
\date{2026-06-20}

\begin{document}
\maketitle

\section*{Status}

This packet gives a strengthened partial result for the natural question of
Haller--Kuuseok--Poldvere~\cite{HKP25}.  The source paper proves that
\(L_1(\mu,X)\) has property \(\CWOS\) when \(X\) is weakly uniformly rotund
(\(w\)UR), and asks whether \(w\)UR can be replaced by the formally weaker
assumption that \(X\) itself has property \(\CWOS\).

The result here proves the source theorem for every strictly convex Banach
space \(X\), with no Radon--Nikodym assumption on \(X^*\).  This is still not
the full \(X\) has \(\CWOS\) problem, because there are \(\CWOS\) spaces that
are not strictly convex, and it does not settle the injective tensor-product
question.

\section*{Source Problem Evidence}

The source question appears in the introduction of~\cite{HKP25}, page 3 of the
arXiv PDF:

\begin{center}
\includegraphics[width=0.95\linewidth]{figures/open_problem_crop.png}
\end{center}

\section*{Transcription}

After proving the \(w\)UR-valued \(L_1\)-theorem and countable
\(\ell_1\)-sum stability for \(\CWOS\), the authors ask whether the theorem
remains true when \(X\) is assumed to be any Banach space with property
\(\CWOS\), and whether the injective tensor product of two \(\CWOS\) spaces
must have property \(\CWOS\).

\section*{Definitions}

A Banach space \(X\) has property \(\CWOS\) if, for every finite convex
combination \(C=\sum_{k=1}^n\lambda_k U_k\) of relatively weakly open subsets
of \(B_X\), every point of \(C\cap S_X\) is an interior point of \(C\) in the
relative weak topology of \(B_X\).

\section*{Proof Intuition}

The source proof uses \(w\)UR to turn almost norm-attainment of local Bochner
averages into weak closeness of the two local summands.  Strict convexity gives
a different mechanism.  If \(z=\lambda x+\nu y\in S_{L_1(\mu,X)}\), then
equality in the \(L_1\)-triangle inequality holds almost everywhere.  In a
strictly convex target this forces \(x(\omega),y(\omega),z(\omega)\) to lie on
the same positive ray almost everywhere.  After partitioning the range of the
unit direction \(e(\omega)=z(\omega)/\|z(\omega)\|\) into small-diameter
pieces, the local averages of \(x,y,z\) are norm-close in the only way needed
by the decomposition argument.  General weak functionals on \(L_1(\mu,X)\) are
handled by composing the finitely many functionals defining the weak
neighbourhoods with the coordinate multipliers \(f\mapsto \one_D f\).

\begin{lemma}[Directional averaging]\label{lem:direction}
Let \(X\) be a Banach space, let \(D\) be measurable, and suppose
\[
  x=r e,\qquad y=s e,\qquad z=q e
\]
on \(D\), where \(e:D\to S_X\), \(r,s,q\ge0\), \(q=\lambda r+\nu s\), and
\(\lambda,\nu>0\), \(\lambda+\nu=1\).  Put
\[
  \alpha=\int_D r\,d\mu,\qquad
  \beta=\int_D s\,d\mu,\qquad
  \gamma=\int_D q\,d\mu.
\]
If \(\gamma>0\) and \(\operatorname{diam} e(D)\le \delta\), then
\[
  \left\|\frac{\alpha}{\gamma}z-x\right\|_{L_1(D)}
  \le 2\delta\alpha,
  \qquad
  \left\|\frac{\beta}{\gamma}z-y\right\|_{L_1(D)}
  \le 2\delta\beta .
\]
\end{lemma}

\begin{proof}
Choose \(e_0\in e(D)\).  Since
\[
  \int_D\left(\frac{\alpha}{\gamma}q-r\right)\,d\mu=0,
\]
we have
\[
  \frac{\alpha}{\gamma}z-x
  =\left(\frac{\alpha}{\gamma}q-r\right)(e-e_0).
\]
Also \(q\ge \lambda r\), so \(\alpha/\gamma\le1/\lambda\), and
\[
  \int_D\left|\frac{\alpha}{\gamma}q-r\right|\,d\mu
  \le \frac{\alpha}{\gamma}\int_D q\,d\mu+\int_D r\,d\mu
  =2\alpha.
\]
The first estimate follows.  The second is identical.
\end{proof}

\begin{theorem}\label{thm:strict}
Let \(\mu\) be a non-zero \(\sigma\)-finite countably additive non-negative
measure and let \(X\) be a strictly convex Banach space.  Then
\[
  L_1(\mu,X)
\]
has property \(\CWOS\).
\end{theorem}

\begin{proof}
By the countable \(\ell_1\)-sum stability theorem of~\cite[Proposition
1.2]{HKP25}, it is enough to prove the result when \(\mu\) is finite.
Indeed, a \(\sigma\)-finite measure space decomposes into countably many
finite-measure pieces, and \(L_1(\mu,X)\) is the corresponding countable
\(\ell_1\)-sum.

Set \(Z=L_1(\mu,X)\), and let \(U,V\) be relatively weakly open subsets of
\(B_Z\).  Let \(x\in U\), \(y\in V\), and \(\lambda,\nu>0\),
\(\lambda+\nu=1\), satisfy
\[
  z=\lambda x+\nu y\in S_Z.
\]
By the reduction lemma of~\cite[Lemma 2.1]{HKP25}, we may assume that \(x\)
and \(y\) have the same support up to null sets: wherever one is non-zero, so
is the other.  It suffices to find a relative weak neighbourhood \(W\) of
\(z\) in \(B_Z\) such that \(W\subset\lambda U+\nu V\).

Choose a finite set \(F\subset S_{Z^*}\) and \(\eta>0\) so small that
\[
  U_0=\{u\in B_Z: |f(u-x)|<\eta\text{ for every }f\in F\}\subset U
\]
and
\[
  V_0=\{v\in B_Z: |f(v-y)|<\eta\text{ for every }f\in F\}\subset V.
\]
All constants below are chosen sufficiently small relative to \(\eta\),
\(\lambda\), \(\nu\), and \(F\).

Since \(\|z\|=\lambda\|x\|+\nu\|y\|=1\), equality holds in the pointwise
triangle inequality almost everywhere.  Strict convexity and the common
support assumption imply that there are measurable functions
\[
  e:\Omega\to S_X,\qquad r,s,q:\Omega\to[0,\infty)
\]
such that
\[
  x=r e,\qquad y=s e,\qquad z=q e,\qquad q=\lambda r+\nu s
\]
almost everywhere on \(\{z\ne0\}\), and all four functions vanish off that set.

Choose \(\beta>0\) and a measurable set \(G\subset\{z\ne0\}\) so that
\[
  r/q\ge\beta,\qquad s/q\ge\beta\quad\text{on }G,
\]
and
\[
  \int_{\Omega\setminus G}(r+s+q)\,d\mu
\]
is as small as needed.  This is possible because \(r/q\) and \(s/q\) are
positive almost everywhere on \(\{z\ne0\}\).  Since \(e\) is strongly
measurable, partition a further large subset of \(G\) into finitely many
measurable sets \(D_2,\dots,D_n\) on which \(\operatorname{diam} e(D_i)\le
\delta\), leaving a bad set \(D_1\) with
\(\int_{D_1}(r+s+q)\,d\mu\) as small as needed.  Put also
\(D_0=\Omega\setminus(D_1\cup\dots\cup D_n)\), which is null for \(z\).  We
choose the partition so that
\[
  \Gamma:=\sum_{i=2}^n\gamma_i>0,
\]
where
\[
  \alpha_i=\int_{D_i}r\,d\mu,\quad
  \beta_i=\int_{D_i}s\,d\mu,\quad
  \gamma_i=\int_{D_i}q\,d\mu .
\]
On the good pieces \(i\ge2\), one has
\[
  \alpha_i/\gamma_i\ge\beta,\qquad \beta_i/\gamma_i\ge\beta.
\]

Using~\cite[Lemma 2.2(b)]{HKP25}, choose a relative weak neighbourhood
\(W_{\rm mass}\) of \(z\) in \(B_Z\) such that, for \(w\in W_{\rm mass}\),
the quantities
\[
  \eta_i(w)=\|\one_{D_i}w\|-\gamma_i,\qquad i=0,\dots,n,
\]
satisfy \(\sum_{i=0}^n|\eta_i(w)|<d\), and also
\(\|\one_{D_1}w\|\) is as small as needed.  Here \(d>0\) is chosen so small
that
\[
  \frac{2d}{\Gamma}\max\{1/\lambda,1/\nu\}<\beta/2.
\]
For each \(f\in F\) and each \(i=1,\dots,n\), the map
\[
  w\mapsto f(\one_{D_i}w)
\]
is a continuous linear functional on \(Z\).  Intersect \(W_{\rm mass}\) with
conditions making every
\[
  |f(\one_{D_i}(w-z))|
\]
smaller than a chosen tiny constant.  Call the resulting relative weak
neighbourhood \(W\).

Fix \(w\in W\), and write \(w_i=\one_{D_i}w\).  For \(i\) with
\(\gamma_i=0\), put \(u_i=v_i=w_i\).  For \(i=1\) with \(\gamma_1>0\), put
\[
  u_1=(\alpha_1/\gamma_1)w_1,\qquad
  v_1=(\beta_1/\gamma_1)w_1.
\]
For the good pieces \(i=2,\dots,n\), define
\[
  \kappa=\sum_{\gamma_i>0}\frac{\beta_i-\alpha_i}{\gamma_i}\eta_i(w),
  \qquad
  \rho=\sum_{i=2}^n\|w_i\|.
\]
The choice of \(d\) gives
\[
  |\kappa|/\rho<\beta/2.
\]
Hence the following coefficients are non-negative:
\[
  \alpha_i/\gamma_i+\nu\kappa/\rho,\qquad
  \beta_i/\gamma_i-\lambda\kappa/\rho\qquad(i=2,\dots,n).
\]
Set
\[
  u_i=\left(\alpha_i/\gamma_i+\nu\kappa/\rho\right)w_i,\qquad
  v_i=\left(\beta_i/\gamma_i-\lambda\kappa/\rho\right)w_i
  \quad (i=2,\dots,n),
\]
and finally \(u=\sum_i u_i\), \(v=\sum_i v_i\).
Then \(\lambda u+\nu v=w\).

The usual \(\ell_1\)-bookkeeping, identical to the computation in the proof of
\cite[Theorem 1.1]{HKP25}, gives \(\|u\|\le1\) and \(\|v\|\le1\).  Indeed,
after summing over all pieces with \(\gamma_i>0\), the term containing
\(\kappa\) changes
\[
  \sum_i(\alpha_i/\gamma_i)\|w_i\|
\]
into
\[
  \sum_i\alpha_i+\sum_i\eta_i(w),
\]
and the remaining \(\gamma_i=0\) pieces satisfy \(x=y=z=0\).  Since
\(\sum_i\eta_i(w)=\|w\|-\|z\|\le0\), we get \(\|u\|\le1\), and the proof for
\(v\) is symmetric.

It remains to prove \(u\in U_0\) and \(v\in V_0\).  Let \(f\in F\).  The bad
piece \(D_1\) contributes only a norm-small term, because
\(\int_{D_1}(r+s+q)\,d\mu\) and \(\|w_1\|\) are small.  On a good piece
\(D_i\), \(i\ge2\), Lemma~\ref{lem:direction} gives
\[
  \left\|(\alpha_i/\gamma_i)z_i-x_i\right\|_1\le2\delta\alpha_i.
\]
Therefore
\[
  u_i-x_i
  =(\alpha_i/\gamma_i)(w_i-z_i)
    +\bigl((\alpha_i/\gamma_i)z_i-x_i\bigr)
    +(\nu\kappa/\rho)w_i .
\]
The first term is small under \(f\) by the definition of \(W\), since
\(\alpha_i/\gamma_i\le1/\lambda\).  The second term is norm-small by the
directional averaging lemma.  The third term has total norm at most
\(|\kappa|\), which is small.  Summing over the finitely many pieces gives
\[
  |f(u-x)|<\eta.
\]
Thus \(u\in U_0\).  The same argument, using the second estimate in
Lemma~\ref{lem:direction}, gives \(v\in V_0\).  Hence \(w=\lambda u+\nu v\in
\lambda U+\nu V\).

Since every \(w\in W\) belongs to \(\lambda U+\nu V\), \(z\) is a relative
weak interior point of \(\lambda U+\nu V\).  This proves \(\CWOS\) for
\(L_1(\mu,X)\).
\end{proof}

\begin{remark}
The proof deliberately avoids identifying \(L_1(\mu,X)^*\) with
\(L_\infty(\mu,X^*)\).  Instead, it localizes the finitely many functionals
that define the weak neighbourhoods by composing them with the bounded
coordinate multipliers \(f\mapsto\one_D f\).  Thus no Radon--Nikodym property
of \(X^*\) is used.
\end{remark}

\section*{Verification Notes}

The proof is non-computational.  The points for review are:
\begin{itemize}
  \item the strict-convexity step forcing pointwise positive collinearity;
  \item the finite partition of the direction map \(e=z/\|z\|\);
  \item the use of localized functionals \(f\circ(\one_D\cdot)\) instead of
        an \(L_\infty(\mu,X^*)\)-representation;
  \item the \(\ell_1\)-norm bookkeeping copied from the source proof.
\end{itemize}

\section*{Limitations}

This is still a partial result.  It proves the \(L_1(\mu,X)\) theorem for the
large natural class of strictly convex target spaces, but it does not cover
all \(\CWOS\) spaces.  In non-strictly-convex \(\CWOS\) spaces, equality in
the pointwise triangle inequality no longer forces a common direction, and the
argument above no longer supplies norm-small local errors.  The injective
tensor-product question remains open here as well.

\section*{Novelty and Literature Check}

I searched the run indexes and local corpus for arXiv:2503.09182, \(\CWOS\),
\(L_1(\mu,X)\), strictly convex targets, and injective tensor variants.  The
existing run packet for this paper proves the arbitrary \(\ell_1\)-sum/atomic
case, and the older arXiv:1806.10693 packet treats a separate finite-dimensional
tensor subcase.  Neither contains the strictly convex \(L_1(\mu,X)\) argument.

Bounded web searches for ``\(L_1(\mu,X)\) strictly convex CWO-S'', ``strictly
convex property CWO-S Bochner'', and close variants found the source paper and
nearby CWO literature, but no exact statement of the theorem above.  Novelty
confidence is moderate pending expert review.

\section*{Human Review Recommendation}

Review the proof as a claimed strict-convex upgrade of the source theorem, not
as a full solution of the \(X\) has \(\CWOS\) question.  The most important
line to audit is the final estimate of \(f(u-x)\): it combines localized weak
control of \(\one_{D_i}(w-z)\), norm-small fixed directional errors, and the
small global correction \(\kappa\).

\begin{thebibliography}{9}

\bibitem{HKP25}
R. Haller, P. Kuuseok, and M. Poldvere,
\emph{Norm-one points in convex combinations of relatively weakly open subsets
of the unit ball in the spaces \(L_1(\mu,X)\)},
arXiv:2503.09182, 2025.
\url{https://arxiv.org/abs/2503.09182}

\bibitem{ABHLP}
T. A. Abrahamsen, J. Becerra Guerrero, R. Haller, V. Lima, and M. Poldvere,
\emph{Banach spaces where convex combinations of relatively weakly open
subsets of the unit ball are relatively weakly open},
Studia Math. \textbf{250} (2020), no. 3, 297--320.
\url{https://arxiv.org/abs/1806.10693}

\end{thebibliography}

\end{document}
